\chapter{Fase 2 -- Grade vertical nativa}
\label{cap:fase2}

\section{Objetivo}

A Fase~2 calcula a geometria vertical completa da malha de destino --
altura física das interfaces entre níveis de modelo ($z_{grid}$), a
métrica vertical $zz$, coeficientes de interpolação/derivação vertical, e
termos de métrica para advecção de terceira ordem -- a partir apenas do
terreno e da conectividade da malha de destino, \textbf{sem nenhuma
dependência do dado meteorológico}. Diferentemente das demais fases, essa
etapa é puramente geométrica e determinística, o que permitiu uma
estratégia de validação particularmente forte: comparar o resultado
diretamente contra a grade vertical de um \texttt{init.nc} real de
produção, sem depender de nenhuma rodada nova.

As fórmulas desta seção foram extraídas literalmente do bloco
\texttt{config\_tc\_vertical\_grid} da rotina \texttt{init\_atm\_case\_gfs}
do código-fonte de referência do MPAS-Model (versão de produção
\texttt{mpas-bundle-3.0.2}, confirmada como o fonte exato que gerou os
binários usados no cluster de produção via inspeção do
\texttt{CMakeCache.txt} do sistema de build) -- não uma reformulação, mas
uma tradução linha a linha, removendo apenas a infraestrutura de troca de
halo entre blocos MPI (desnecessária nesta implementação, que roda serial
sobre a malha inteira).

\section{Perfil vertical de referência $z_w(k)$}
\label{sec:zw}

O perfil de altura das interfaces de nível na configuração
``2014 TC experiments'' (identificada, por inspeção do \textit{namelist}
real de produção, como o branch de código efetivamente usado em todas as
rodadas deste projeto -- não a fórmula genérica mais simples de versões
anteriores do MPAS) é definido por trechos, com constantes calibradas
para $n_{VertLevels}=55$:
\begin{equation}
  z_w(k) =
  \begin{cases}
    \left[ a_s \dfrac{k-1}{n_1} + c_2\left(\dfrac{(k-1)\Delta z}{z_t \zeta_l}\right)^{\!2} - c_3\left(\dfrac{(k-1)\Delta z}{z_t \zeta_l}\right)^{\!3} \right] z_t, & \dfrac{k-1}{n_1} < \zeta_l \\[1.2em]
    \left[ z_l + a_t\left(\dfrac{k-1}{n_1} - \zeta_l\right) \right] z_t, & \dfrac{k-1}{n_1} \ge \zeta_l
  \end{cases}
  \label{eq:zw}
\end{equation}
com $n_1 = n_{VertLevels}$, $\Delta z = z_t/n_1$, $z_l = 1 - a_t(1-\zeta_l)$,
$c_2 = 3(1-a_t) + 2(a_t-a_s)\zeta_l$, $c_3 = 2(1-a_t) + (a_t-a_s)\zeta_l$,
e as constantes calibradas $a_s = 0{,}075$, $a_t = 1{,}70$,
$\zeta_l = 0{,}75$ (para $n_{VertLevels} \ge 55$; valores distintos,
$a_t=1{,}23$ e $\zeta_l=0{,}31$, aplicam-se ao caso legado de 41 níveis,
não usado neste projeto). O topo do domínio $z_t$ é o parâmetro de
configuração \texttt{config\_ztop} (\SI{30000}{\metre} em todos os casos
de produção deste projeto).

\section{Coordenada híbrida e o achado do bloqueio}
\label{sec:ah_hibrido}

A função de atenuação $A(\zeta)$ introduzida na
\cref{eq:coord_hibrida_teoria} determina a rapidez com que a influência
do terreno é removida das superfícies de coordenada com a altura. O
código-fonte de referência oferece dois ramos:
\begin{equation}
  A(k) =
  \begin{cases}
    \cos^6\!\left( \dfrac{\pi}{2} \dfrac{z_w(k)}{z_h} \right), & z_w(k) < z_h \quad \text{(coordenada híbrida)} \\[0.6em]
    0, & z_w(k) \ge z_h \\[0.6em]
    1 - \dfrac{z_w(k)}{z_t}, & \text{(coordenada \textit{terrain-following} básica)}
  \end{cases}
  \label{eq:ah}
\end{equation}
onde $z_h$ é a altura de transição (parâmetro \texttt{config\_hybrid\_top\_z}).

\begin{caixaachado}
\label{achado:hybrid}
Este ponto foi o principal bloqueio numérico enfrentado no
desenvolvimento desta fase, e vale ser documentado em detalhe pelo valor
metodológico da investigação. A suposição inicial, baseada na ausência
de qualquer menção a \texttt{config\_hybrid\_coordinate} nos arquivos
\texttt{namelist.init\_atmosphere} reais de produção, foi de que essa
opção usava seu valor padrão \texttt{.false.} -- isto é, o ramo
\textit{terrain-following} básico (segunda linha da
\cref{eq:ah}). Sob essa suposição, a grade vertical calculada
divergia da referência real em até $\sim\SI{2}{\kilo\metre}$ em níveis
intermediários -- um erro fisicamente inaceitável, da mesma classe de
gravidade dos bugs catastróficos já documentados no relatório principal
do projeto (NaN em \texttt{PSFC}).

Uma primeira hipótese, investigada extensivamente, foi de que o processo
de suavização iterativa (\cref{sec:suavizacao}) seria genuinamente não
reprodutível bit a bit entre uma reimplementação serial e a rodada real
distribuída por MPI, dada a natureza do critério de aceitação por limiar
usado nesse algoritmo (pequenas diferenças de arredondamento podem levar
células individuais a aceitar ou rejeitar a atualização de forma
diferente, compondo-se ao longo de dezenas de passadas). Essa hipótese
foi \emph{descartada} experimentalmente: reimplementar toda a cadeia de
suavização em precisão simples (\texttt{real*4}), para casar exatamente
com a precisão do binário de produção (confirmada pelo log real,
``Default real precision: single''), não produziu nenhuma melhora
mensurável -- os valores de erro permaneceram idênticos, dígito a dígito.

A causa raiz real só foi encontrada ao baixar o arquivo
\texttt{Registry.xml} do \texttt{core\_init\_atmosphere} diretamente do
repositório de referência do MPAS-Model: a opção
\texttt{config\_hybrid\_coordinate} tem \texttt{default\_value="true"}
(com \texttt{config\_hybrid\_top\_z} padrão \SI{30000}{\metre}) --
confirmado de forma ainda mais direta ao acessar o código-fonte exato
usado para compilar o binário de produção (\texttt{mpas-bundle-3.0.2}),
onde a coordenada híbrida está, na verdade, \emph{fixada no
código-fonte} (\texttt{hybrid = .true.}, não sequer exposta como opção de
\textit{namelist} nessa versão). Ou seja: a ausência da chave no
\textit{namelist} real não significava ``desativado'', significava
``usando o padrão -- que é ativado''. Corrigindo a fórmula de $A(k)$ para
o primeiro ramo da \cref{eq:ah}, o erro máximo na grade vertical caiu de
$\sim\SI{1103}{\metre}$ para $\SI{171}{\metre}$, e o erro médio de
$\SI{27,4}{\metre}$ para $\SI{0,054}{\metre}$ -- uma redução de
aproximadamente 500 vezes, com o resíduo restante concentrado quase
inteiramente no anel de células de fronteira da malha regional
(\cref{sec:validacao_fase2}), não mais distribuído por toda a coluna
vertical.

A lição metodológica: uma revisão linha a linha do bloco de código
suspeito (o laço de suavização) não encontrou nenhuma diferença de
lógica porque a causa real estava numa dependência \emph{upstream} desse
bloco (o cálculo de $A(k)$, calculado dezenas de linhas antes) -- ao
investigar uma divergência numérica que resiste a uma revisão de código
aparentemente correta, é necessário validar também as dependências do
trecho suspeito, não apenas o próprio trecho.
\end{caixaachado}

\section{Suavização do terreno e da grade vertical}
\label{sec:suavizacao}

Antes de aplicar as fórmulas anteriores, o campo de terreno bruto passa
por uma suavização espacial de quarta ordem (filtro de Shapiro,
$n_{smterrain}$ passadas, tipicamente $1$), com coeficiente $+0{,}216$
seguido de $-0{,}216$ em cada passada -- validado como essencialmente
exato ($\text{erro médio}\approx\SI{0,15}{\metre}$) contra a referência
real.

Mais crítico é o algoritmo de suavização \emph{iterativa por nível} da
superfície de terreno $h_x(k,\cdot)$, que progressivamente remove
estrutura de pequena escala do terreno conforme a altura aumenta -- a
técnica descrita originalmente por Klemp
\cite{klemp2011terrainfollowing} (Eqs.~5--9 daquele trabalho, confirmadas
como estruturalmente idênticas ao código real do MPAS-A nesta
investigação). Para cada nível $k=2,\dots,n_{VertLevels}-1$, executam-se
$n_{sm}+k$ passadas de um filtro Laplaciano local ponderado pela
resolução da malha,
\begin{equation}
  h_x^{(n)}(i) = h_x^{(n-1)}(i) + s(i)\!\!\sum_{j \in \mathcal{N}(i)}\!\! \frac{\ell_{ij}}{d_{ij}}\Big[h_x^{(n-1)}(j) - h_x^{(n-1)}(i)\Big],
  \label{eq:suavizacao_hx}
\end{equation}
onde $\mathcal{N}(i)$ é o conjunto de células vizinhas de $i$,
$\ell_{ij}$/$d_{ij}$ são o comprimento da aresta e a distância entre
centros de célula, e $s(i)$ é um fator de suavização local escalado pela
resolução nominal da malha e pela altura do nível corrente. A atualização
só é \emph{aceita} se o espaçamento vertical resultante não cair abaixo
de uma fração mínima configurável do espaçamento nominal
($\delta z_{min} > \texttt{config\_dzmin} \cdot \Delta z_w(k)$) --
critério que existe justamente para impedir camadas degeneradas
(espaçamento nulo ou negativo) perto de relevo íngreme.

\section{Campos derivados e coeficientes numéricos}

A partir de $z_w(k)$ e $A(k)$, calculam-se diretamente (sem depender de
mais nenhuma suavização): $\Delta z_w(k) = z_w(k{+}1)-z_w(k)$,
$rdz_w(k)=1/\Delta z_w(k)$, $\Delta z_u(k)=\tfrac12[\Delta z_w(k)+\Delta z_w(k{-}1)]$,
$rdz_u(k)=1/\Delta z_u(k)$, e os coeficientes de projeção de interface
$f_{zm}(k)$, $f_{zp}(k)$ (interpolação linear entre camadas adjacentes),
além dos três coeficientes de extrapolação de contorno de topo
$cf_1,cf_2,cf_3$. Esse conjunto de campos, por depender apenas de
\texttt{config\_ztop}/$n_{VertLevels}$/\texttt{config\_interface\_projection}
-- sem depender da geometria da malha em si -- foi validado como
essencialmente exato (diferenças da ordem de $10^{-6}$ a $10^{-9}$,
atribuíveis a arredondamento de armazenamento em precisão simples) contra
a referência real.

Finalmente, a altura física das interfaces e a métrica vertical são
\begin{equation}
  z_{grid}(k, i) = z_w(k) + A(k)\, h_x(k, i),
  \qquad
  zz(k, i) = \frac{z_w(k{+}1) - z_w(k)}{z_{grid}(k{+}1, i) - z_{grid}(k, i)},
  \label{eq:zgrid_zz}
\end{equation}
e os termos de métrica horizontal $zx_u$ (usados na projeção do gradiente
de altura ao longo das arestas) são calculados por diferença finita
central de $z_{grid}$ entre as duas células vizinhas de cada aresta,
normalizada pela distância entre centros de célula ($dc_{Edge}$).

\subsection{Termos de métrica para advecção de terceira ordem ($zb$, $zb^3$)}

Um segundo achado de implementação, corrigindo uma suposição do plano
original, foi que os campos $zb$/$zb^3$ -- usados no cálculo do vento
vertical diagnóstico da Fase~4 (\cref{cap:fase4}) para o esquema de
advecção de terceira ordem de $\theta$ -- \textbf{não} vêm prontos no
arquivo de malha estática, ao contrário do que se assumia inicialmente;
eles dependem de $z_{grid}$ (produto desta própria fase) e precisam ser
recalculados após ela, usando os coeficientes de segunda derivada
horizontal $\texttt{deriv\_two}$ (esses sim já presentes no arquivo de
malha estática). Uma nova subrotina (\texttt{compute\_zb}, incluída no
mesmo módulo desta fase) foi escrita para esse fim, extraída literalmente
do bloco ``For $z$-metric term in omega equation'' do código de
referência.

\section{Validação numérica}
\label{sec:validacao_fase2}

Após a correção do \cref{achado:hybrid}, a grade vertical calculada foi
comparada campo a campo contra o $z_{grid}$/$zz$ real de um
\texttt{init.nc} de produção (malha \texttt{SouthAmerica}, 17064 células,
56 interfaces). O erro máximo caiu para $\SI{171}{\metre}$ e o erro médio
para $\SI{0,054}{\metre}$. A análise da distribuição espacial do resíduo
confirmou que ele está \textbf{concentrado quase inteiramente} nas sete
camadas do anel de fronteira da malha regional
(\texttt{bdyMaskCell}~$\in\{1,\dots,7\}$, cerca de 2600 das 17064
células): nas células estritamente interiores
(\texttt{bdyMaskCell}$=0$, 82\% da malha), o erro máximo cai para
$\SI{10,9}{\metre}$ -- fisicamente irrelevante. Essa correlação espacial
tem explicação física conhecida: a rotina de referência aplica um passo
de \emph{mistura de terreno de fronteira}
(\texttt{config\_blend\_bdy\_terrain}) \emph{antes} de chamar a geração
da grade vertical, misturando a altura de terreno das células de
fronteira com o terreno do dado de fundo -- passo ainda não implementado
nesta rota (discutido como limitação conhecida no
\cref{cap:discussao}).
